\documentclass[11pt,a4paper]{article}
\usepackage[T1]{fontenc}
\usepackage[utf8]{inputenc}
\usepackage[british]{babel}
\usepackage{amsmath,amssymb,booktabs,tabularx,graphicx,float}
\usepackage{geometry,microtype,csquotes}
\usepackage[font=small,labelfont=bf]{caption}
\usepackage[backend=biber,style=authoryear,maxcitenames=2,maxbibnames=6,
            doi=true,isbn=false,url=false]{biblatex}
\usepackage[hidelinks]{hyperref}
\geometry{margin=22mm}
\addbibresource{references.bib}
\graphicspath{{../results/technical-report/figures/}}
\setlength{\parindent}{0pt}
\setlength{\parskip}{0.45em}
\renewcommand{\arraystretch}{1.12}
\newcommand{\reporttable}[1]{%
  \input{../results/technical-report/tables/#1.tex}}
\reporttable{results-macros}
\hypersetup{pdftitle={Antibody duration: methodology and initial results},
  pdfsubject={Prior comparisons and repeated single-participant fits}}

\title{\textbf{Antibody duration: methodology and initial results}\\[0.3em]
  \large Prior comparisons and repeated single-participant fits}
\author{}
\date{15 October 2026}

\begin{document}
\maketitle
\vspace{-2em}
\section{Models, estimands and calibration}
Time is in days and titres are relative to mean convalescent neutralisation.
For each fixed threshold $c_e$, $e\in\{\mathrm{sym},\mathrm{sev}\}$, define
\begin{equation}
  T_e^*=\inf\{t\geq0:c(t)\leq c_e\},\qquad \inf\varnothing=\infty.
\end{equation}
The target is individual duration. A trajectory initially at or below a
threshold has $T_e^*=0$; the same posterior curve determines both durations,
with $T_{\mathrm{sym}}^*\leq T_{\mathrm{sev}}^*$.

The fitted families are
\begin{align}
  \text{Exponential:}\quad c(t)&=c_0e^{-kt}, & k&>0;\label{eq:exp}\\
  \text{Biphasic:}\quad c(t)&=c_0\{(1-f)e^{-k_1t}+fe^{-k_2t}\},
    & k_1&>k_2>0;\label{eq:bip}\\
  \text{Power law:}\quad c(t)&=c_0(1+\alpha t)^{-\beta},
    & \alpha,\beta&>0;\label{eq:power}\\
  \text{Plateau:}\quad c(t)&=A_\infty+(c_0-A_\infty)e^{-kt},
    & k&>0,\quad 0<A_\infty<c_0.\label{eq:plateau}
\end{align}
For the biphasic curve, $f$ is the initial titre fraction contributed by
the slow component. Its half-lives and equal-contribution time satisfy
\begin{equation}
  k_r=\frac{\log2}{h_r},\quad
  f=\frac{1}{1+e^{(k_1-k_2)t_s}},\quad
  t_s=\frac{\log\{(1-f)/f\}}{k_1-k_2}.
\end{equation}
Uncensored observations follow
\begin{equation}
  \log y_j=\log c(t_j)+\epsilon_j,\qquad
  \epsilon_j\overset{\mathrm{iid}}\sim N(0,\sigma_{\log}^2).
\end{equation}
There are no random effects or censored observations in this pilot.

\begin{table}[H]
\centering\small
\caption{Generating parameters and fixed analysis anchors. Titres and
thresholds use the convalescent-relative scale.}
\reporttable{calibration}
\label{tab:calibration}
\end{table}
The biphasic calibration is from \textcite{hogan2023}; the thresholds
are fixed anchors for 50\% symptomatic/severe protection from
\textcite{khoury2021}. Threshold uncertainty is not propagated here.

For completeness, the proposed later mechanistic generator is
\begin{equation}
  \dot P_s=-\delta_sP_s,\qquad \dot P_l=-\delta_lP_l,\qquad
  \dot c=\rho_sP_s+\rho_lP_l-\delta_cc.
\end{equation}
It is not used in the present simulation or fits.

\section{Prior formulations}
Let $t_a=90$ and define
\begin{equation}
  D_a=-\log\{c(t_a)/c_0\}>0,\qquad R_a=e^{-D_a}.
\end{equation}
All four specifications share the appropriate default/broader column in
Table~\ref{tab:common}. $N(\mu,s^2)$ denotes a normal distribution with
standard deviation $s$. All free variables are independent within a model;
no additional priors are assigned to derived parameters.

\begin{table}[H]
\centering\small
\caption{Common priors in both formulations.}
\label{tab:common}
\reporttable{common-priors}
\end{table}

\subsection{Early-only formulation}
Every family has $c(0)=c_0$ and $c(t_a)=c_0R_a$.

\paragraph{Exponential.} Set $k=D_a/t_a$ in Equation~\eqref{eq:exp}.

\paragraph{Biphasic.} With free $u_f\in(0,1)$ and $z_q\in\mathbb R$, set
\begin{equation}
  f=u_f/2,\qquad q=1+e^{z_q},\qquad
  (1-f)e^{-qst_a}+fe^{-st_a}=R_a,\qquad k_2=s,\quad k_1=qs.
\end{equation}
The equation has a unique $s>0$. Here $f<0.5$ and $t_s>0$.

\paragraph{Power law.} With free $\beta>0$, set
\begin{equation}
  \alpha=(e^{D_a/\beta}-1)/t_a.
\end{equation}

\paragraph{Plateau.} With free $u_\infty\in(0,1)$, set
\begin{equation}
  A_\infty=u_\infty c_0R_a,\qquad r=A_\infty/c_0,\qquad
  k=-t_a^{-1}\log\!\left(\frac{R_a-r}{1-r}\right).
\end{equation}

\begin{table}[H]
\centering\small
\caption{Additional priors for the early-only default and broader
specifications. The rate-ratio centre is the value used in the code,
$16.6\simeq581/35$.}
\label{tab:early}
\reporttable{early-shape-priors}
\end{table}

\subsection{Two-anchor formulation}
For biphasic, power-law and plateau models, let $t_b=365$, $m=t_b/t_a$,
and introduce $\rho\in(0,1)$:
\begin{equation}
  \rho=\frac{D_b-D_a}{(m-1)D_a},\qquad
  D_b=D_a\{1+(m-1)\rho\},\qquad R_b=e^{-D_b}.
\end{equation}
Thus $R_a^m<R_b<R_a$ and all three families share the joint prior on
$c_0,c(90),c(365)$. The exponential retains $k=D_a/t_a$ and its
early-only prior; it has no free $\rho$.

\paragraph{Power law.} Solve for $v>0$ and set
\begin{equation}
  \frac{\log\{1+m(e^v-1)\}}{v}=\frac{D_b}{D_a},\qquad
  \beta=D_a/v,\qquad \alpha=(e^v-1)/t_a.
\end{equation}

\paragraph{Biphasic.} With the sole additional free variable
$u_s\in(0,1)$, set
\begin{equation}
  k_2=u_s\rho D_a/t_a,\quad x=e^{-k_2t_a},\quad
  f=\frac{R_a-y}{x-y},\quad
  (1-f)y^m+fx^m=R_b,\quad k_1=-\log(y)/t_a.
  \label{eq:anchor-mix}
\end{equation}
The root $y\in(0,R_a)$ is unique. Here $0<f<1$; $f$ and $t_s$ are
derived, and $f>0.5$ (hence $t_s<0$) is permitted.

\paragraph{Plateau.} Use Equation~\eqref{eq:anchor-mix} with $x=1$,
then set $A_\infty=c_0f$ and $k=-\log(y)/t_a$.

\begin{table}[H]
\centering\small
\caption{Additional priors for the two-anchor default and broader
specifications. These replace Table~\ref{tab:early}.}
\label{tab:anchors}
\reporttable{two-anchor-priors}
\end{table}
The reference curve gives $D_a\simeq1.16$, $\rho\simeq0.25$ and
$u_s\simeq0.37$. Priors are fixed across reviews. The day-365 anchor
is a latent parameter, not an added observation.

\section{Pilot design and computation}
One participant is simulated from the biphasic calibration in
Table~\ref{tab:calibration}. The \ObservationCount{} measurements are at
$t_j=\operatorname{round}(365.25j/12)$, $j=0,\ldots,24$; the same noisy
history is truncated at 91, 183, 274, 365, 548 and 730 days (labelled
3, 6, 9, 12, 18 and 24 months). All four models are fitted under each
of the four prior specifications. True individual crossing times are
\TruthSym{} days (symptomatic) and \TruthSev{} days (severe).

JAGS 5.0.0 uses four chains, initially 2,000 adaptation, 2,000 warm-up
and 4,000 retained iterations per chain. Up to two retries multiply draws
by four and warm-up by two. Require rank-normalised split
$\widehat R<1.01$ and bulk/applicable tail ESS $\geq400$ for free
parameters and duration functions. Diagnose $T^*/(365.25+T^*)$ (one at
infinity) and non-constant non-crossing indicators. Constant
outcome/quantile indicators have no ESS. Seeds derive from the design,
replicate, model, prior specification and chain.

Intervals below are equal-tailed 90\% posterior intervals. Non-crossing
draws remain at $\infty$. Prior curves use 4,000 draws per family and
specification; common RNG streams match the shared anchors. Extreme
power-law prior durations are summarised on log time. Subsequent-visit
predictions are evaluated only after fitting, using
\begin{equation}
  \ell_j=\log\left\{\frac1S\sum_{s=1}^S
    p\bigl(\log y_j\mid\theta^{(s)}\bigr)\right\},\qquad
  \overline\ell=\frac1{n_{\mathrm{future}}}\sum_j\ell_j.
\end{equation}
The future visits lie between the current and next review. Pointwise
PSIS-LOO is descriptive; Pareto-$k>0.7$ flags unreliability and it does
not validate long-term extrapolation.

\section{Initial results}
\paragraph{Computation.} Of 96 labelled model--prior--review combinations,
\FitCount{} completed and all \ConvergedCount{} passed the declared checks,
across \AttemptsCount{} sampling attempts including the numerical failure
below. The two formulations
give identical exponential posterior draws within each prior width.
The broader two-anchor power law at three months failed to emit a valid
duration result because some finite crossing times exceeded floating-point
range. Its raw sampler output is retained; it is marked as missing rather
than recoded as non-crossing. Only \ReliableLooCount{} of \FitCount{}
completed fits had reliable pointwise PSIS-LOO diagnostics.

\paragraph{Prior implications.} Matching the first-year anchors leaves
substantial tail disagreement (Figures~\ref{fig:priors} and
\ref{fig:non-cross}). The two-anchor power-law severe-duration prior has
an approximate 95th percentile of $1.0\times10^{11}$ days (default) or
$3.6\times10^{26}$ days (broader). These exploratory priors require
further tail calibration before a main simulation study.

\paragraph{Posterior duration.} The exponential estimates the severe
crossing too early under every specification, with intervals excluding
the true crossing (Figure~\ref{fig:duration}). At 24 months its default estimate is
\EarlyDefaultExp{} days, below the true \TruthSev{} days. Under the
early-only default prior the biphasic estimate is \EarlyDefaultBip{};
under the two-anchor default it is \AnchorDefaultBip{}.
All four biphasic intervals contain the true severe crossing. The four
power-law intervals at 24 months exclude it from above. The plateau
assigns probability one, to reported precision, to no severe crossing
at 24 months. The higher-threshold estimates are much more similar among
the flexible models after longer follow-up.

Early fits can resemble one another while their
tails differ (Figure~\ref{fig:posterior}). Subsequent-visit log scores
(Figure~\ref{fig:future}) assess a shorter extrapolation than the severe
threshold. These pilot results concern one noisy individual history; interval
containment here is not an estimate of empirical coverage.

\clearpage
\begin{figure}[p]
\centering
\includegraphics[width=\textwidth]{prior-trajectories.pdf}
\caption{\textbf{Prior latent trajectories for all four specifications.}
Lines show medians and shading 90\% pointwise intervals. Rows identify the
prior formulation and width; columns identify the model. Dashed lines
are the two thresholds; dotted lines mark days 90 and 365. The
exponential shares only the initial/day-90 priors. The two-anchor
flexible families additionally share the day-365 prior.}
\label{fig:priors}
\end{figure}

\begin{figure}[p]
\centering
\includegraphics[width=\textwidth]{posterior-durations.pdf}
\caption{\textbf{Individual threshold-crossing estimates.} Points are
posterior medians and bars are 90\% intervals; horizontal dashed lines
show generating truth. The display is capped at 20 years: ``Inf'' denotes
an infinite upper limit, ``$>$'' a finite upper limit beyond the cap, and
an open triangle at the cap an infinite median. The cross near the bottom
marks the failed broader two-anchor power-law fit at three months. The
two exponential rows of each width are identical.}
\label{fig:duration}
\end{figure}

\begin{figure}[p]
\centering
\includegraphics[width=\textwidth]{posterior-trajectories.pdf}
\caption{\textbf{Posterior latent trajectories as follow-up accumulates.}
Lines are medians and shading 90\% pointwise intervals. Black dashed
curves are the generating trajectory; black points are observations
available at that review. Dotted lines show the thresholds. Columns
select four reviews; rows show all four prior specifications. The
three-month broader two-anchor power-law curve is omitted because its
duration output failed.}
\label{fig:posterior}
\end{figure}

\begin{figure}[p]
\centering
\includegraphics[width=\textwidth]{non-crossing.pdf}
\caption{\textbf{Probability of never crossing each threshold.}
Open diamonds show prior probabilities; filled points show posterior
probabilities at successive reviews. All positive-rate exponential,
biphasic and power-law draws eventually cross; the plateau permits a
genuine non-crossing outcome. The failed broader two-anchor power-law
fit at three months has no posterior point. These probabilities retain
the infinite-duration mass omitted by a finite-only summary.}
\label{fig:non-cross}
\end{figure}

\begin{figure}[p]
\centering
\includegraphics[width=\textwidth]{future-prediction.pdf}
\caption{\textbf{Prediction of withheld subsequent visits.} Each point is
the mean log predictive density of observed log titres between the
indicated review and the next review; higher is better. The four panels
use identical observations and evaluation windows. No score is available
after the final 24-month review. The failed three-month broader
two-anchor power-law fit is omitted. These scores concern near-term
prediction, not validation of the severe crossing time.}
\label{fig:future}
\end{figure}

\clearpage
\IfFileExists{../results/technical-report/tables/repeated-macros.tex}{%
  % Included by technical-report.tex once the completed study is assembled.
\reporttable{repeated-macros}
\section{Repeated single-participant simulation}

\subsection{Design and scores}
Generate \RepeatN{} independent noisy histories of one participant using
the biphasic truth and assay noise of Table~\ref{tab:calibration} and
the pilot's monthly visits.
Fit all four models at six reviews with the early-only default priors
(Tables~\ref{tab:common} and~\ref{tab:early}). Histories are shared across
models and nested reviews. This is a fixed-truth measurement-noise study
without between-individual heterogeneity.

For posterior median $m_r$ and central 90\% interval $[L_r,U_r]$, estimate
coverage and bias among the $R_{\mathrm{ok}}$ converged replicates as
\begin{equation}
  \widehat C=R_{\mathrm{ok}}^{-1}\sum_r
    \mathbf1\{L_r\leq T^*\leq U_r\},\qquad
  \widehat B=R_{\mathrm{ok}}^{-1}\sum_r(m_r-T^*).
\end{equation}
Precision is summarised by the median interval width and unbounded-interval
fraction. Intervals with either endpoint infinite count as unbounded;
infinite medians remain in the bias calculation. Coverage bars are Wilson
95\% intervals. For finite replicate-level means, Monte Carlo SE is the
sample SD divided by $\sqrt{R_{\mathrm{ok}}}$; bars show mean
$\pm1.96$ SE. At the planned \RepeatN{} replicates, the worst-case
coverage SE is \RepeatCoverageMaxSe{} (\RepeatCoverageNominalSe{} at
90\% coverage). These SEs describe variation across histories; within-fit
MCMC uncertainty is controlled by the convergence checks.

Ordinary day-scale CRPS diverges in the presence of non-crossing mass.
Use the one-to-one transformation $g(t)=t/(365.25+t)$, with
$g(\infty)=1$, to score the entire extended duration distribution:
\begin{equation}
  \operatorname{CRPS}_g(F,T^*)=
  \int_0^1\{F_g(z)-\mathbf1[g(T^*)\leq z]\}^2\,dz.
\end{equation}
Also report restricted CRPS in days,
\begin{equation}
  \operatorname{CRPS}_H(F,T^*)=
  \int_0^H\{F(t)-\mathbf1[T^*\leq t]\}^2\,dt,
  \qquad H=3652.5\text{ days}.
\end{equation}
This equals CRPS for $\min(T^*,H)$ and scores only the distribution up to
ten years; five- and twenty-year horizons are sensitivity checks. Lower
scores are better \parencite{gneiting2007}. All empirical CDF integrals are
computed exactly between sorted draws. The scoring transformations do not
alter the saved crossing times or non-crossing probabilities. Bias is also
reported on the bounded $g$ scale, where it remains finite.

\subsection{Results}
Of \RepeatPlanned{} fits, \RepeatCompleted{} completed and
\RepeatConverged{} converged across \RepeatAttempts{} sampling attempts.
Numerical/engine failures: \RepeatErrors{}; unresolved fits:
\RepeatUnconverged{}. Model/review denominators are
\RepeatMinDenominator{}--\RepeatMaxDenominator{}. Failure-inclusive
coverage bounds assign unsuccessful fits first a miss and then a hit.
All attempts and outcomes remain in the saved outputs.

For severe duration, exponential coverage is \RepeatExpEarlyCoverage{}
at three months and \RepeatExpLateCoverage{} at 24 months; biphasic
coverage is \RepeatBipEarlyCoverage{} and \RepeatBipLateCoverage{},
respectively (Figure~\ref{fig:repeat-coverage}). Power-law and plateau
coverage at 24 months is \RepeatPowerLateCoverage{} and
\RepeatPlateauLateCoverage{}. Mean
posterior-median bias is \RepeatExpLateBias{} days for the exponential,
\RepeatBipLateBias{} for biphasic and \RepeatPowerLateBias{} for power
law. Plateau bias is \RepeatPlateauLateBias{} when any replicate has an
infinite median; this is not replaced by a finite-only average.

At 24 months, mean severe-duration $\operatorname{CRPS}_g$ is
\RepeatExpLateCrps{} (exponential), \RepeatBipLateCrps{} (biphasic),
\RepeatPowerLateCrps{} (power law) and \RepeatPlateauLateCrps{} (plateau).
Figures~\ref{fig:repeat-precision}--\ref{fig:repeat-crps} distinguish
precision from accuracy: a narrow interval may concentrate on the wrong
outcome, including infinity. The restricted day score reverses the
exponential/power-law ordering; both scores favour biphasic at 24 months.
Figure~\ref{fig:repeat-horizon} checks the restricted-score horizon.
Results are specific to this fixed reference
curve, schedule, noise and prior set.

\clearpage
\begin{figure}[p]
\centering
\includegraphics[width=\textwidth]{repeated-coverage.pdf}
\caption{\textbf{Repeated-simulation coverage.} Coverage of nominal 90\%
individual-duration intervals across \RepeatN{} histories per model/review.
Whiskers are Wilson 95\% intervals among converged fits; the dashed line
is nominal coverage. Thick translucent segments give all-planned
worst-case coverage bounds for unsuccessful fits. The two endpoint panels
use the same histories.}
\label{fig:repeat-coverage}
\end{figure}

\begin{figure}[p]
\centering
\includegraphics[width=\textwidth]{repeated-width.pdf}
\includegraphics[width=\textwidth]{repeated-unbounded.pdf}
\caption{\textbf{Precision and unbounded intervals.} Top: median width
of central 90\% intervals, retaining unbounded intervals as infinite
width. Display is capped at twenty years; ``Inf'' is an infinite median
width and ``$>$'' a finite median width above the cap. Bottom: fraction of
intervals with an infinite endpoint, with Wilson 95\% intervals. A small
finite-only width would conceal these non-crossing outcomes.}
\label{fig:repeat-precision}
\end{figure}

\begin{figure}[p]
\centering
\includegraphics[width=\textwidth]{repeated-bias.pdf}
\includegraphics[width=\textwidth]{repeated-bias-bounded.pdf}
\caption{\textbf{Bias of the posterior median.} Top: mean error in days,
retaining infinite medians. The display clips values and error bars at
$\pm3652.5$ days; ``Inf'' marks infinite bias and ``$>$'' a finite mean
above the upper cap. Bottom: mean error after applying $g$ to both the
posterior median and the true duration. Whiskers show $\pm1.96$ Monte
Carlo SE where finite; the dashed line is zero bias.}
\label{fig:repeat-bias}
\end{figure}

\begin{figure}[p]
\centering
\includegraphics[width=\textwidth]{repeated-crps.pdf}
\includegraphics[width=\textwidth]{repeated-crps-restricted.pdf}
\caption{\textbf{Distributional accuracy.} Top: mean CRPS on the bounded
$g$ scale, scoring the full distribution including infinity. Bottom:
mean restricted CRPS through ten years, in days. Lower values are better.
Whiskers show $\pm1.96$ Monte Carlo SE; lower error limits are truncated
at zero. Ordinary day-scale CRPS is not used to compare models assigning
non-crossing mass.}
\label{fig:repeat-crps}
\end{figure}

\begin{figure}[p]
\centering
\includegraphics[width=\textwidth]{repeated-horizon.pdf}
\caption{\textbf{Sensitivity to the restricted-CRPS horizon.} Mean
day-scale CRPS restricted to five, ten and twenty years. Each row uses
the same fitted posterior distributions and true crossing times; only the
score integration horizon changes. Scores do not assess timing beyond
the stated horizon.}
\label{fig:repeat-horizon}
\end{figure}
\clearpage}{}
\section*{Reproduction}
From the repository root, run the fitting studies and report assemblers:
\begin{verbatim}
Rscript scripts/01-prior-pilot.R
Rscript scripts/01-prior-pilot.R --two-anchor
Rscript scripts/03-technical-report.R
Rscript scripts/04-repeated-study.R
Rscript scripts/05-repeated-report.R
cd notes
latexmk -pdf -outdir=../results/technical-report technical-report.tex
\end{verbatim}
The editable main source is \path{notes/technical-report.tex}.
The repeated-results section is \path{notes/repeated-results.tex}.
The assemblers validate the saved simulation, fit and score artefacts and write vector
figures, input tables, an input-file checksum manifest and session
information to \texttt{results/technical-report/}. Once fits are saved,
only scripts 03 and 05 and the compiler are needed to rebuild the report.
Priors and calibration follow the handout;
this report introduces no additional model assumptions.

\printbibliography
\end{document}
